% !TeX root = ../papers/nc-p3.tex
\section{一つの二十点配置への射影と平方根の衝突}
\label{sec:single-configuration-projection}
定理\ref{thm:full-fifteen-component}の十五次元の有理商を$M$とする。
第一位置の二十点の射影同値類を取る有理写像を$\pi$とする。

\begin{theorem}\label{thm:single-configuration-projection}
$\pi$は像へ生成的有限で、その関数体次数は偶数である。
第一関係六平面を固定座標で記録した像の次元は45、第一位置の点配置のみを
固定座標で記録した像の次元は30である。$t$および一般点で、四位置の各座標
四空間は自分の平方を固定したGrassmannスキームの孤立した被約点である。
$r$の偶位置ではこの局所ファイバーは長さ2であり、奇位置では被約点である。
\end{theorem}
\begin{proof}
第一切断の左核の基底を$L$とすると、第一関係空間を一定にする接方向は
$L^{\mathsf t}T=0$である。これと第二切断のSchur微分の共通核は$r,t$で
文字次数順に$(7,8,8,8)$次元となる。第三・第四位置の基底変更の軌道は
$\mathrm{GL}_4^2$の逆スカラー一組を除いて31次元で、常にこの共通核に入る。
従って一般にも31次元であり、第一辺の像の次元は$76-31=45$である。

次に任意の第一六平面の基底$r_1,\ldots,r_6$を動かす。点対$(a,b)$で
$E_a=(a^{\mathsf t}r_s)_s$は階数3で、第一点$a$を固定して第二点を動かせる
一次条件は$(a^{\mathsf t}\dot r_s b)_s\in\operatorname{im}E_a$である。
番号付き二十点がエタールに動くので、全二十対について課せば十分でもある。
各対につき三式を並べた60行96列の階数は両例で45である。下界は素数12619,2657で
得られ、上界は六関係の基底変更36と第二位置の射影変換15が核に入ることによる。
後者は先の点グラフと射影枠から独立に作用する。従ってGrassmann多様体の
接空間で第一配置を固定する方向は全て第二位置の基底変更であり、先の45次元像
にもこれらの軌道が含まれる。その第一配置への像は$45-15=30$次元、
第一位置の射影変換を除けば$30-15=15$で、$\pi$の生成的有限性を得る。
全階数の検査は\path{computations/verify_single_configuration_projection.py}にある。

奇二位置の交換は四窓を逆に読む操作で、左右AS条件を保存する。
第一位置の配置は逆向きのグラフの同じ端点となる。この操作は$M$を保つ。
実際、四元群の規格化前には$(M_h,N_h)$を$(N_h^{\mathsf t},M_h^{\mathsf t})$に
移すので、非零な薄いブロック、可逆な交差・組内ブロックを保つ。
薄い積からの再規格化は対角変換である。$z=\pm1$が一般に生じないことも、
$M'_2=N_2^{\mathsf t},\ M'_3=N_3^{\mathsf t}$の薄い係数の組内積比から従う。
より具体的には$\lambda=1$の片側対称な断面で、偶交換と奇交換の合成は
対置行列を転置し、薄い外積の二ベクトルを交換する。このとき薄い係数
$u_i^h v_k^h$の規格化前の$v_i^h=u_i^{h+1}$により両色が交換され、
$U,V$の各二列の順序が反転するので積比$z$は$1/z$に移る。
従って一般の像は定理\ref{thm:klein-all-pairings}の非退化チャートにある。
五次元の像は四つの例外表示（各々印なし次元は高々4）に入り得ず、
楕円曲線族の閉包に戻る。この族が$t$で滑らかな唯一の全成分に入るため
逆転はその成分も保つ。一般点で第二と第四の配置は射影非同値なので、
この対合は関数体に非自明に作用する。$\pi$は対合に不変であり、その次数は偶数となる。

最後に$D=k^{20}$で一位置の座標行列の列空間を$A$とする。
$B^2\subset A^2$となる$B\in\Gr(4,D)$を考え、$A$で独立な四行を固定する。
残り十六行の四係数がGrassmann座標になる。変分$S$の一次式は、全ての
$i\le j$について$A_i S_j+A_j S_i\bmod A^2$である。$t$の四位置と$r$の
二つの奇位置で階数64を得る。$r$の偶位置では還元後の階数は63で、
核の一方向$S$の二次項$(S_i S_j\bmod A^2)_{ij}$を付け加えると64となる。

$r$の偶位置の標数零での接次元が少なくとも1である理由も必要である。
二十点を番号付けした成分上、第一・第三位置の平方は単元倍を除いて等しい。
これは各辺の完全交叉のtraceを比較する定理\ref{thm:finite-square-chart}の証明で、
有限エタールな族上でもtraceとその比は局所切断となる。
必要な行ごとの平方根をエタール局所で選ぶ。$r$では第一・第三の点列が一致し、
行の自明化を合わせると平方を保つ単元は命題\ref{prop:finite-square-ratio}によりスカラーである。従って符号を合わせれば二つの平方根部分空間は一致する。
一般点では両配置が射影非同値なので二切断は相異なる。相対Grassmannスキームが
この点で不分岐なら対角がその近傍で開となり、二切断は一般点でも一致してしまう。
従って標数零でも接次元は1である。単元な63次小行列で核の生成元を
素数12619の剰余へ整合的に取れるため、上の二次障害も標数零で非零となる。
六十三変数を形式的に消去すると残りは一変数で、局所イデアルはその二乗で生成される。
これで重なった平方根の局所長さを得る。
\end{proof}

\begin{corollary}
十五次元成分の幾何学的一般点では、第一配置だけの射影安定群も自明である。
\end{corollary}
\begin{proof}
$t$で、上の二段の接計算は、$Z$の接空間を基底変更の方向で割った空間から
第一配置の接空間を射影変換の方向で割った空間への微分が単射であることを示す。
実際、第一配置を射影的に固定する接方向から第一位置の基底変換を引き、
次に第二位置の基底変換を引けば第一辺が一定になり、最後は第三・第四位置の軌道に入る。
一般の配置の安定群の元が与える置換を、有限エタールな番号付け上で$t$に特殊化する。
枠と特異行列の除外により、それは$t$の第一配置の四元群の元$g$となり、
枠からの行列は近傍で正則に動く。この$g$は$t$の四窓にも作用する。
従って固定等式を微分すると、全ての窓の商接方向$v$について
$\mathrm d\pi(gv-v)=0$が配置の商接空間で成立する（ここでは群軌道を引いた
接空間同士の写像を用いる）。単射性と全非自明文字が現れることから$g=1$である。
番号付け上で異なる置換は特殊点でも異なるため結論を得る。
\end{proof}

\paragraph{一つの配置からの有限候補。}
一般ファイバーの復元には、第一座標の四空間$A$から$Q=A^2\subset D$を作り、
$Q$と標準内積による$Q^\perp$に対して平方が入る四空間を求めればよい。
第一辺の完全交叉のtraceは各点で非零なので、第二座標の行単元を調整すれば
その平方は$Q^\perp$になる。第三座標は同様に$Q$の平方根である。
一般ファイバーの実際の解は上の孤立被約点に入るので、この二つの
Grassmannスキームの零次元成分を列挙すれば必ず候補に含まれる。
全階数の有理点では、四辺とも、第一空間自身を第三空間として再使用する候補は三次交叉で落ちる。実際、各六関係を列とする$R$について$R\otimes I_4$と$I_4\otimes R^{\mathrm{op}}$を並べると有理数上で階数48となり、交叉次元は必要な4に達しない。最大階数は一般点にも保たれる。残る各候補の隣接積の核から二窓を作り、系\ref{cor:two-window-extension}で
延長と同型を判定する。次節では生成的有限写像の正確な次数を決める。特殊配置での正則性判定はなお必要である。

従って復元次数$e$には、粗いが一様な上界
\[
 2\le e\le 2^{45}(2^{45}-1),\qquad 2\mid e
\]
がある。実際、一つの平方根スキームの孤立点は有限個なので、全てを一つの一般の
Grassmannアフィン座標に入れられる。六十四変数において平方条件は高々二次であり、
孤立解に対するBézout不等式からその数は$2^{64}$以下となる（正次元成分がある場合も
孤立点の局所重複度の和に同じ上界が成立する）。実際の四空間では平方を安定化する
対角単元はスカラーのみであるから、座標の符号倍は全体符号を除いた$2^{19}$個の
相異なる孤立解を与える。各符号軌道は同じ積の核を与える。よって実際の候補は
各側高々$2^{45}$通りで、第三位置では既知の第一空間の符号軌道を除ける。
二窓延長の一意性により積が上界となる。
